\section{三个 Product-Critical Systems：Indexing、Inference 与 Apply}

课程先把系统拆成 indexing、model inference 和 product/apply layer。它们都服务同一个 editor，却有完全不同的资源曲线：indexing 偏异步与 data-intensive；inference 受 GPU capacity、batching 和 latency 约束；apply layer 需要把生成结果快速而安全地映射回用户正在编辑的文件。

\subsection{不同 workload 需要不同 SLO}

本节把上一段的 workload 差异转成可观测合同。读图时先为 indexing、inference、apply 和 data plane 分别定义“成功、降级和失败”，再选择 freshness、recall、queue time、edit success 或 job correctness 等指标；如果所有路径只共享一个 availability 数字，事故时就无法判断应当牺牲哪类工作来保护交互主链。

\lecturefigure{02-three-systems.png}{Cursor 的规模来自 indexing、inference、product/apply 与 data plane 四类负载。}{本地字幕 02:20--05:20；概念重绘。}

\begin{knowledgebox}{读图：先给每个系统定义失败}
Indexing 的失败可能是 stale、missing 或低 recall；inference 的失败可能是 timeout、rate limit 或低质量；apply 的失败可能是 patch 冲突、位置错误或 editor freeze；data plane 的失败可能是 job 丢失、重复执行或 billing 不可解释。只有先分类型，报警和降级才不会互相污染。
\end{knowledgebox}

\begin{importantbox}{SLO 不能只写“99.9\% available”}
\begin{itemize}
  \item Index freshness：workspace change 到可检索的延迟；
  \item Retrieval quality：在给定 latency 和 filter 下的 recall/utility；
  \item Inference：time-to-first-token、tokens/s、error 与 queue time；
  \item Apply：edit success、conflict、validation latency 与 undo safety。
\end{itemize}
\end{importantbox}

\subsection{Apply model：把“想改什么”与“如何落盘”分开}

规划模型擅长理解意图、选择文件和生成修改；specialized apply model 则处理位置匹配、局部重写和高速 token application。把两者分开，是因为最强模型未必在低延迟、确定性 patch application 上最优。系统需要为“思考”和“写入 editor state”设计不同接口。

\teachervoice{课堂反复强调产品层并不是模型输出后的薄 UI。真正让工具“感觉好用”的，往往是 apply、streaming、context preparation 和 editor integration 这些不容易在 benchmark 中出现的系统工作。}

\subsection{本章小结}

AI coding 平台不是一个后端服务，而是多个 SLO 和数据形状不同的系统组合。产品价值来自这些系统在一次交互中的协同，而不是某个单点模型指标。

